\subsection{Critical-locus variance refinements}\label{sec:fqe-renyi-compatibility}
The universal bound (QK) suffices for the anchor argument. We now give its sharper
critical-locus refinements: the fixed-variance profile on $P=1$, used to classify equality
in (QK), and the analytic proof of the sharp bound~\eqref{eq:fqe-quadratic-variance} on every critical state.
Put $p_j=\operatorname{tr}\rho_{1j}^2$ and $U=\sum_jp_j^2$.

\begin{proposition}[Sharp purity-variance tradeoff at minimal purity]
\label{prop:fqe-critical-variance}
Suppose a pure four-qubit state has $P=1$. Set
\[
V=\sum_{j=2}^4(p_j-1/3)^2,\qquad s=(24V)^{1/4}.
\]
Then $0\le s\le1$ and the sharp bound at each fixed $V$ is
\[
\boxed{\quad F\ge\frac32s^2-s^3.\quad}
\tag{RF}\label{eq:fqe-critical-variance}
\]
For $a=(1+2s)/3$, $b=(1-s)/3$, equality holds exactly when, up to ordering of
cuts and eigenvalues, the three spectra are
\[
\begin{split}
&\left(\frac12,\frac a2,\frac b2,\frac b2\right),\\
&\left(\frac{1+b+2\sqrt{ab}}4,\frac{1+b-2\sqrt{ab}}4,
\frac{1-b}4,\frac{1-b}4\right)\quad\text{on each of the other two cuts}.
\end{split}
\tag{RF-eq}\label{eq:fqe-variance-spectra}
\]
These spectra are attained for every $s\in[0,1]$.
In particular, $T+U\le3/4$, equivalently $F\ge12V$; equality in this weaker
bound occurs only at the HS spectra ($s=0$) or cluster spectra ($s=1$).
\end{proposition}
\noindent[Proved under the hypothesis $P=1$.]
\begin{proof}
By Lemma~\ref{lem:fqe-purity}, the state is critical with $\tau=0$. As in the proof of
Proposition~\ref{prop:fqe-critical-envelope}, its Bloch correlation matrix is then $\Theta=\mathbf m\mathbf n^T$ with unit
$\mathbf m,\mathbf n\in\mathbb R^3$ (Lemma~\ref{lem:fqe-critical-correlation}(2) at $\kappa=0$), and the cut $c\in\mathbb Z/3$ has the
four eigenvalues $(1+\sigma\cdot g^{(c)})/4$, $g^{(c)}_a=m_{a+c}n_a$, as $\sigma$ runs over the sign triples with
$\sigma_1\sigma_2\sigma_3=-1$ (Lemma~\ref{lem:fqe-p1-cuts}). Put $A_a=m_a^2$, $B_a=n_a^2$ and $k=-m_1m_2m_3n_1n_2n_3$.
The three vectors $g^{(c)}$ have squared norms
\[
u_c=\sum_aA_{a+c}B_a,\qquad \sum_cu_c=1,
\]
and the same coordinate product $-k$. The sign sums $\sum_\sigma\sigma\cdot g=0$, $\sum_\sigma(\sigma\cdot g)^2=4|g|^2$ and
$\sum_\sigma(\sigma\cdot g)^3=-24g_1g_2g_3$ give, for the corresponding cut,
\[
p_j=\frac{1+u_c}{4},\qquad
\operatorname{tr}\rho_{1j}^3=\frac{1+3u_c+6k}{16}.
\]
Set $X=1-3e_2(A)$ and $Y=1-3e_2(B)$, where $e_2(A)=\sum_{i<j}A_iA_j$.
Then $0\le X,Y\le1$, and expanding the three cyclic sums gives
\[
\sum_c u_c^2=\frac{1+2XY}{3}.
\]
We sharpen the elementary Schur estimate by keeping $X$ and $Y$.
For any nonnegative triple $A$ of sum one and $X=1-3e_2(A)$,
\[
27e_3(A)\le h(X),\qquad
h(X)=1-3X+2X^{3/2}=(1-\sqrt X)^2(1+2\sqrt X).
\tag{RP}\label{eq:fqe-triple-product}
\]
This is the cubic discriminant identity
\[
 4X^3-\bigl(27e_3(A)-1+3X\bigr)^2
   =27\prod_{i<j}(A_i-A_j)^2\ge0.
\]
Taking the upper square-root bound proves~\eqref{eq:fqe-triple-product}.
Equality requires two equal coordinates, and the upper sign selects
$((1+2r)/3,(1-r)/3,(1-r)/3)$, $r=\sqrt X$, up to permutation.
This includes the uniform triple and coordinate-vector endpoints.

The transfer from $X,Y$ to their geometric mean is algebraic as well.
Set $g(t)=(1-t)^2(1+2t)$, $\xi=X^{1/4}$, $\zeta=Y^{1/4}$,
$q=\xi\zeta$, $d=(\xi-\zeta)^2$, and
$w=(1-\xi^2)(1-\zeta^2)$. Direct factorisation gives
\begin{equation}\label{eq:fqe-product-transfer-remainder}
 \begin{split}
 h(\sqrt{XY})^2-h(X)h(Y)
 &=g(\xi\zeta)^2-g(\xi^2)g(\zeta^2)\\
 &=d\bigl\{w[6q(q+2)+2d]+d(1+2q)^2\bigr\}\ge0.
 \end{split}
\end{equation}
Every factor is nonnegative for $0\le X,Y\le1$, and the difference
is strictly positive when $X\ne Y$. Consequently
\[
27k\le27\sqrt{e_3(A)e_3(B)}
\le\sqrt{h(X)h(Y)}\le h(\sqrt{XY}).
\]
The moment formulas above give
\[
V=\frac{XY}{24},\qquad F=\frac12-\frac{27}{2}k.
\]
Thus $s=(XY)^{1/4}$ and $F\ge(1-h(s^2))/2=3s^2/2-s^3$.
For $0<s<1$, equality requires $X=Y=s^2$, $k>0$, and both $A,B$ to be
permutations of $(a,b,b)$. The shift $c$ that aligns their distinguished
coordinates gives $|g^{(c)}_a|=(a,b,b)$ up to order; the other two shifts give
$(\sqrt{ab},\sqrt{ab},b)$. As every $g^{(c)}$ has negative coordinate product,
all sign choices give the same spectral multisets, and the eigenvalue formula
proves \eqref{eq:fqe-variance-spectra}.
At $s=0$, equality forces $k=1/27$ and hence both triples uniform; at $s=1$
both are coordinate vectors. These give the stated endpoint spectra.

For attainability, take $\mathbf m=-\mathbf n=(\sqrt a,\sqrt b,\sqrt b)$, so $A=B=(a,b,b)$ and $k=ab^2=e_3(A)$.
The product auxiliary vector with Bloch vectors $\mathbf m,\mathbf n$ has $\Theta=\mathbf m\mathbf n^T$ and $\tau=0$, and it is
$\sum_kv_km_k$ for a critical $G(v)$, as at the end of the proof of Proposition~\ref{prop:fqe-critical-envelope}.
Finally,
\[
\frac32s^2-s^3-12V=\frac{s^2}{2}(1-s)(3+s)\ge0,
\]
with equality exactly at $s=0,1$. This proves the weaker bound and its
exhaustive equality statement as a corollary.
\end{proof}

\subsection{A joint-cut proof on the critical locus}
\label{sec:fqe-critical-quadratic-variance}

\begin{proposition}[Sharp quadratic variance on critical states]
\label{prop:fqe-critical-quadratic-variance}
Every critical unit four-qubit state satisfies
\[
 F^2+3P(P-1)\ge6V_P.
 \tag{CQ}\label{eq:fqe-critical-quadratic-variance}
\]
Equality holds at HS spectra and along the full cluster-to-Bell family
\[
 \bigl((1+\kappa)/2,(1-\kappa)/2,0,0\bigr),\qquad
 (1/4)^4,\qquad(1/4)^4,\qquad 0\le\kappa\le1,
\]
up to ordering of cuts and eigenvalues, and only there.
In particular, the defect in \eqref{eq:fqe-quadratic-variance-proved} is at least
$3\tau(1-\tau)/4$ on the critical locus. The proof uses no variance certificate. At $P=1$ the sharper
profile~\eqref{eq:fqe-critical-variance} of Proposition~\ref{prop:fqe-critical-variance} holds.
\end{proposition}
\begin{proof}
Use the auxiliary correlation matrix of Proposition~\ref{prop:fqe-critical-envelope}:
\[
 \Theta=U\operatorname{diag}(\kappa,-\kappa,1)V^T,\qquad U,V\in SO(3),
 \qquad \tau=\kappa^2,\qquad P=1+\kappa^2/2.
\]
Write $r=Ue_3$, $s=Ve_3$, $A_i=r_i^2$, $B_i=s_i^2$, and
$X=1-3e_2(A)$, $Y=1-3e_2(B)$. Thus $0\le X,Y\le1$.
We retain these two directional parameters in both the permanent and purity estimates.

\emph{The permanent bound with directions retained.}
Put $z=\sqrt{XY}$ and
\[
 b_\kappa(x)=\kappa^2+\frac23\kappa(1-\kappa)^2(1-x)
             +\frac29(1-\kappa)^3(1-3x+2x^{3/2}).
\]
The polar-decomposition Cauchy--Schwarz bound from
Proposition~\ref{prop:fqe-critical-envelope}, now using
\eqref{eq:fqe-triple-product} instead of discarding $X,Y$, gives
\[
 |\operatorname{perm}\Theta|
 \le\sqrt{b_\kappa(X)b_\kappa(Y)}\le b_\kappa(z).
\]
For completeness, the last inequality is geometric concavity.
Write $b_\kappa(x)=a-bx+cx^{3/2}$, where
\[
 a=(4\kappa^3+3\kappa^2+2)/9,\quad
 b=2(1-\kappa)^2/3,\quad c=4(1-\kappa)^3/9.
\]
The sign of $d^2\log b_\kappa(e^v)/dv^2$ is the negative of
$4ab-9ac t+bc t^3$, $t=\sqrt x$.
This polynomial decreases on $[0,1]$, since $b\le3a$, and its value at one is
$4\kappa^2(1-\kappa)^2(4\kappa^2+\kappa+1)/9\ge0$.
The cases with zero arguments follow by continuity.
Consequently the exact formula for $F$ gives
\[
 F\ge\kappa^2(1-\kappa)+\frac32(1-\kappa)^2z
                  -(1-\kappa)^3z^{3/2}
 \ge f_\kappa(z),
\]
where, using $\sqrt z\le(1+z)/2$,
\[
 f_\kappa(z)=(1-\kappa)\kappa^2
       +\frac12(1-\kappa)^2(1+2\kappa)z
       +\frac12(1-\kappa)^3z(1-z)\ge0.
\]

\emph{The three purities as one Fourier coefficient.}
Let $\omega=e^{2\pi i/3}$, $D=\operatorname{diag}(1,\omega,\omega^2)$, and
\[
 u_c=\sum_{j=0}^2\Theta_{j+c,j}^2,\qquad
 Z=\sum_c\omega^c u_c
   =\operatorname{tr}(D\Theta\overline D\Theta^T).
\]
Indices are cyclic. The cut formula $p_c=(1+u_c)/4$ implies
$6V_P=|Z|^2/4$.
The symmetric unitary matrices $M=U^TDU$ and $N=V^T\overline D V$
have $|M_{33}|^2=X$, $|N_{33}|^2=Y$.
Their last-row off-diagonal squared norms are $1-X,1-Y$;
their leading $2\times2$ blocks have squared Frobenius norms $1+X,1+Y$.
Expanding $Z$ in these blocks and applying Cauchy--Schwarz gives
\[
 |Z|\le z+2\kappa\sqrt{(1-X)(1-Y)}
                +\kappa^2\sqrt{(1+X)(1+Y)}
       \le B_\kappa(z):=(1-\kappa)^2z+2\kappa+\kappa^2.
\]
To see the last step, fix $z$ and put $S=X+Y\ge2z$.
The last two terms become
$2\kappa\sqrt{1-S+z^2}+\kappa^2\sqrt{1+S+z^2}$,
which decreases with $S$ for $0\le\kappa\le1$; evaluate at $S=2z$.

\emph{An elementary scalar comparison.}
For $0\le\kappa,z\le1$,
\[
 4f_\kappa(z)^2+9\kappa^2-B_\kappa(z)^2
       -3\kappa^2(1-\kappa^2)
   =(1-\kappa)^2(1-z)R_\kappa(z),
\]
where, putting $d=1-\kappa$,
\[
 R_\kappa(z)=-d^4z^3+3d^3(1+\kappa)z^2
                    -4\kappa d^2z+2\kappa^2(1+2\kappa^2).
\]
Since $z^3\le z^2$,
\[
 R_\kappa(z)\ge
 2d^3(1+2\kappa)z^2-4\kappa d^2z+2\kappa^2(1+2\kappa^2)\ge0.
\]
The final quadratic is nonnegative: its discriminant is
$-16\kappa^3d^3(3+2\kappa+4\kappa^2)\le0$ and its leading and constant
coefficients are nonnegative (the endpoints follow directly).
Combining the estimates proves
\[
 4F^2+9\tau-24V_P\ge3\tau(1-\tau),
\]
which is exactly~\eqref{eq:fqe-critical-quadratic-variance}.

For $0<\kappa<1$ the scalar comparison is strict unless $z=1$.
Then $X=Y=1$, and equality in the Fourier triangle bound forces all nonzero
entries of $\Theta$ onto one cyclic diagonal. Their magnitudes are
$1,\kappa,\kappa$, with negative product since $\det\Theta=-\kappa^2$.
The cut eigenvalue formula gives precisely the displayed family.
At $\kappa=0$, equality requires $z=0$ or $z=1$; the latter is its cluster
endpoint, while the former requires $F=0$ and hence HS by the critical
envelope equality case. At $\kappa=1$, equality requires $F=0$ and
$|Z|=3$, forcing the Bell endpoint.
Conversely the family spectra are attained by the correlation matrices
$\operatorname{diag}(\kappa,-\kappa,1)$, and direct substitution gives equality.
\end{proof}


